 % 本文件由 scripts/assemble.py 自动装配，勿手改（改 parts/ 后重跑）
% 第6章 更一般的黑洞

% 第6章 More General Black Holes
\chapter{更一般的黑洞}

\section{黑洞动物园}

% ---- 书页 238 / p253 ----
Birkhoff 定理保证了 Schwarzschild 度规是广义相对论中唯一的球对称真空解。这并不太令人惊讶，因为它让我们想起电磁学中的情形：在无电荷的区域里，唯一的球对称场位形就是 Coulomb 场。走出球对称性之后，可能的引力场就五花八门、不胜枚举。比如对一颗像地球这样的行星，其外部场将依赖于表面上所有各种山脉与山谷的密度和轮廓。我们可以设想把度规按多极矩（multipole moments）展开；要精确描述这个场，就得给定无穷多个系数。

因此，黑洞并不具备这种性质，这一点或许反倒会让人感到几分意外。稳态黑洞解只有寥寥几个，每个解由少数几个参量描述。参量的具体集合取决于我们的理论中包含哪些物质场；如果电磁学是唯一的长程非引力场，我们就有一条\textbf{无毛定理}（no-hair theorem）：

与电磁学相耦合的广义相对论的稳态、渐近平直黑洞解，只要在事件视界之外无奇点，就完全由质量、电荷与磁荷以及角动量这几个参量刻画。

稳态解之所以特别令人感兴趣，是因为我们预期它们是引力坍缩的终态。另一种可能是某种振荡的位形，但随着能量通过引力辐射的发射而损耗，振荡终将被阻尼；实际上，典型的演化会相当迅速地趋于稳态位形。

我们说的是“一条”无毛定理，而不是“那条”无毛定理，因为这一结果不仅依赖于广义相对论，还依赖于我们理论中的物质内容。在粒子物理的标准模型中，电磁学是唯一的长程场，上面的定理适用；但对别的种类的场来说，可能存在其他种类的毛发。\footnote{相关讨论见 M. Heusler, “Stationary Black Holes: Uniqueness and Beyond,” \emph{Living Rev.\ Relativity} \textbf{1} (1998): 6, http://www.livingreviews.org/Articles/Volume1/1998-6heusler/.} 甚至还有人找到了静态（无转动）的、轴对称却不完全球对称的黑洞解的例子。

% ---- 书页 239 / p254 ----
然而核心结论没有改变：黑洞解由极少数参量刻画，而不像比如说一颗行星那样，需要可能无穷多的参量来刻画。

正如我们将在本章末尾以及在第 9 章再次讨论的那样，无毛性质导致了一个令人困惑的局面。在大多数物理理论中，我们希望有一个良定义的初值问题，从而关于任一时刻之态的信息都可以用来预言（或回溯）任何其他时刻的态。这样一来，由运动方程的一个解联系起来的任何两个态，就应当需要同样多的信息来给定。但在 GR 中，事情似乎是这样的：我们可以取一团非常复杂的物质，让它坍缩成黑洞，最终得到的位形却完全由质量、电荷和自旋来描述。在经典 GR 中这或许还不会太困扰我们，因为可以认为这些信息只是藏在事件视界背后，而非真正丢失。但一旦把量子场论考虑进来，我们就会发现黑洞会蒸发并最终消失，信息看起来就真的丢失了。可以想见，导致蒸发的出射霍金辐射或许以某种方式编码了关于原来是用哪个态来造出黑洞的信息，但这一切如何发生完全不清楚。在许多人看来，理解这个“信息丢失佯谬”（information loss paradox）是构建一门合理的量子引力理论的关键一步。

不过在本章中，我们将只限于讨论经典 GR。我们先对黑洞的性质作一些一般性的讨论，尤其是事件视界和 Killing 视界的性质。这个话题可以相当微妙和技术化，我们在这里的方针是力图传达主要思想，而不在定义和定理证明上讲求严格。然后我们讨论带电（Reissner–Nordström）黑洞和转动（Kerr）黑洞所对应的具体解；与我们的一贯做法一致，我们不会仔细梳理构造最大延拓时空所需的那些坐标重新定义，而是径直画出相应的共形图。想进一步了解细节的读者，可以查阅 Townsend 的综述文章\footnote{P.K. Townsend, “Black Holes: Lecture Notes,” http://arxiv.org/abs/gr-qc/9707012.}，以及 Hawking 与 Ellis (1973) 和 Wald (1984) 的书；本章大量取材于这些文献。

\section{事件视界}

黑洞的特点在于：你可以进入它，却永远无法出来。因此，黑洞最重要的特征其实并不是位于中心的奇点，而是位于边界的事件视界。事件视界是一个超曲面，它把能用类时路径与无穷远相连通的时空点，同那些不能连通的时空点分隔开来。要在实践中理解这句话的含义，我们需要更仔细地想一想“无穷远”指的是什么。在广义相对论中，时空的整体结构可以呈现出多种不同的形态，相应地也就有各不相同的无穷远概念。但要思考真实宇宙中的黑洞，我们
% ---- 书页 240 / p255 ----
其实并不真正关心无穷远处发生的事情；我们把无穷远当作“远在黑洞之外”的代名词，并且设想离洞足够远的时空可以用 Minkowski 空间来近似。

\begin{figure}[H]
\centering
\includegraphics[width=0.72\textwidth]{fig_6.1.pdf}
\caption*{图 6.1\quad 渐近平直时空是指：其共形图中的无限远与 Minkowski 时空的无限远一致，包括未来类光无限远 $\mathscr{I}^{+}$、类空无限远 $i^{0}$ 和过去类光无限远 $\mathscr{I}^{-}$。未来事件视界是 $\mathscr{I}^{+}$ 的过去的边界。虚线区域代表时空的其余部分，在不同的例子中它可以呈现多种不同的形式。}
\end{figure}
正如第 5 章提到的，在无穷远处看起来像 Minkowski 时空的时空被称为渐近平直（asymptotically flat）时空。这个概念的含义在一幅诸如图 6.1 那样的共形图中一目了然。从附录 H 中关于 Minkowski 时空共形图的讨论我们知道，共形无限远分为五块：未来与过去类时无限远 $i^{\pm}$、未来与过去类光无限远 $\mathscr{I}^{\pm}$，以及类空无限远 $i^{0}$。如果一个时空（或时空的一个区域）的 $\mathscr{I}^{\pm}$ 和 $i^{0}$ 具有与 Minkowski 时空相同的结构，它就是渐近平直的；类时无限远则不是必需的。这样的时空一般具有图 6.1 所示的形态。

有了这幅图像，我们应该如何理解未来事件视界就清楚了：它是类时曲线无法再逃逸到无穷远的那一侧边界曲面。回忆一个区域的因果过去 $J^{-}$，是指从该区域出发、沿指向过去的类时路径所能到达的所有点的集合；事件视界于是可以等价地定义为 $J^{-}(\mathscr{I}^{+})$——即未来类光无限远的因果过去——的边界。（事件视界其实应当是这个集合的\emph{闭包}（closure）的边界，不过我们在这里不追求严格。）对过去视界也有类似的定义。正如在对 Schwarzschild 作最大延拓的情形中所看到的，一个时空中可以不止有一个渐近平直区域，相应地也就不止有一个事件视界。

从定义可以看出，事件视界是一个类光超曲面。类光超曲面的性质在附录 D 中讨论过；这里我们回顾一下其主要特点。一个超曲面 $\Sigma$ 可以由某个函数 $f(x)$ 通过 $f(x)=$ 常数来定义。梯度 $\partial_\mu f$ 与 $\Sigma$ 正交；若法矢量是类光的，就称这个超曲面为类光超曲面，此时法矢量同时也切于 $\Sigma$。类光超曲面可以看成是一族类光测地线 $x^{\mu}(\lambda)$ 的集合，这些测地线称为该超曲面的生成元（generators）。这些测地线的切矢量 $\xi^{\mu}$ 与法矢量成比例，

% ---- 书页 241 / p256 ----
\begin{equation*}
\xi^\mu = \frac{dx^\mu}{d\lambda} = h(x)g^{\mu\nu}\partial_\nu f,
\tag{6.1}
\end{equation*}
因而也可以充当超曲面的法矢量。我们可以选取函数 $h(x)$，使得这些测地线是仿射参量化的，于是切矢量将满足
\begin{equation*}
\xi_\mu\xi^\mu = 0, \qquad \xi^\mu\nabla_\mu\xi^\nu = 0.
\tag{6.2}
\end{equation*}
对于未来事件视界，其生成元可以在过去终结（例如，当黑洞由恒星坍缩形成时），但总会一直延伸到无穷远的未来（未来与过去互换后亦然）。

由于事件视界是一个整体性的概念，如果有人递给你一个用任意坐标写出的度规，要真正找到它的视界位置可能相当困难。幸运的是，本章我们要打交道的都是相当特殊的度规——稳态、渐近平直，而且含有拓扑为球面的事件视界。在这类时空中，存在一些方便的坐标系，使得辨认事件视界有一种简单办法。对 Schwarzschild 解来说，事件视界是光锥“倾倒”的地方，使得 $r=2GM$ 成为一个类光曲面，而不像在大 $r$ 处那样 $r=$ 常数是类时曲面。光锥倾倒显然是一个依赖坐标的概念（比如在 Kruskal 坐标中就不会发生），但我们关心的那些度规都容许类似的构造。稳态度规有一个渐近类时的 Killing 矢量 $\partial_t$，我们可以调整度规分量使之与时间无关（$\partial_t g_{\mu\nu}=0$）。在 $t=$ 常数的超曲面上，我们可以选取坐标 $(r,\theta,\phi)$，使得无穷远处的度规看起来就像球坐标下的 Minkowski 空间。$r=$ 常数的超曲面在 $r\to\infty$ 处将是拓扑为 $S^{2}\times\mathbf{R}$ 的类时柱面。现在设想我们聪明地选好了坐标，使得当 $r$ 从无穷远逐渐减小时，$r=$ 常数的超曲面始终保持类时，直到某个固定的 $r=r_{\mathrm{H}}$ 处，曲面处处类光。（在不那么聪明的坐标中，$r=$ 常数的超曲面会在某些 $\theta$ 和 $\phi$ 值处变成类光或类空，而在其他值处仍是类时。）这显然就代表一个事件视界，因为跨入 $r<r_{\mathrm{H}}$ 区域的类时路径永远无法逃回无穷远。确定 $r=$ 常数超曲面在何处变为类光很容易：$\partial_\mu r$ 是这种超曲面的一个法 1-形式，其模长为
\begin{equation*}
g^{\mu\nu}(\partial_\mu r)(\partial_\nu r) = g^{rr}.
\tag{6.3}
\end{equation*}
我们要找的正是这个 1-形式的模长为零的地方；于是，在我们所描述的坐标中，事件视界 $r=r_{\mathrm{H}}$ 就是 $g^{rr}$ 由正变负的那个超曲面，
\begin{equation*}
g^{rr}(r_{\mathrm{H}}) = 0.
\tag{6.4}
\end{equation*}

% ---- 书页 242 / p257 ----
这条判据对 Schwarzschild 解显然有效，因为它的 $g^{rr}=1-2GM/r$。后面我们给出 Reissner–Nordström 解和 Kerr 解时，也将采用类似地适配于视界的坐标。

我们之所以把事件视界看得如此重要，是因为在广义相对论中它们几乎是不可避免的。这一结论由两个有趣的结果串联而来：奇点几乎不可避免，而且奇点藏在事件视界背后。当然，两个结果都在各自的一组适当假设下才成立；构造没有奇点的时空并不难（Minkowski 空间就是一例），找到没有视界的奇点也不算难（下面讨论带电黑洞时就会看到）。但我们相信，“通有”（generic）的解都会有藏在视界背后的奇点。

奇点的普遍存在由 Hawking 和 Penrose 的\textbf{奇点定理}（singularity theorems）保证。在这些定理被证明之前，人们还可以指望，坍缩到 Schwarzschild 奇点只是球对称性的一个人为假象，典型的几何将保持无奇点（例如在 Newton 引力中就是这样）。但 Hawking–Penrose 定理表明，坍缩一旦达到某个程度，向奇点的演化就不可避免。我们知道存在奇点的途径是测地不完备（geodesic incompleteness）——存在某条测地线，它无法在流形内延伸，却在仿射参量的有限值处终结。而我们知道坍缩已经达到有去无回之点的标志，是一个\textbf{陷俘面}（trapped surface）的出现。要理解什么是陷俘面，先想象 Minkowski 空间中的一个二维球面，它嵌在一个等时切片里，取为到原点径向距离固定的一些点。如果我们追踪从这个空间球面射向时空的类光线，一族（指向内部）将描绘出一组不断收缩的球面，而另一族（指向外部）将描绘出一组不断胀大的球面。但对 Schwarzschild 几何中半径固定为 $r<2GM$ 的球面来说，情况就不是这样了；在事件视界之内，从这个球面发出的两组类光线都会演化到更小的 $r$ 值（因为 $r$ 是类时坐标），从而演化到更小的面积 $4\pi r^{2}$。这就是陷俘面的含义：一个紧致、类空的二维子流形，其性质是在该子流形上处处，指向未来的出射光线沿两个方向都\emph{汇聚}。（“汇聚”的正式定义是：膨胀 $\theta$——正如附录 F 中关于测地线汇的讨论所描述的——为负。）

有了这些定义，我们就可以给出一个奇点定理的例子。

令 $M$ 是一个流形，带有通有的度规 $g_{\mu\nu}$，满足 Einstein 方程并施加强能量条件。如果 $M$ 中存在一个陷俘面，则必定要么存在一条闭合类时曲线，要么存在一个奇点（表现为某条不完备的类时或类光测地线）。

% ---- 书页 243 / p258 ----
这里所谓“通有度规”，指的是对类时和类光测地线都满足\textbf{普适条件}（generic condition）。对类时测地线，普适条件是说：每条以 $U^{\mu}$ 为切矢量的测地线，都必须至少有一个点使得 $R_{\alpha\beta\gamma\delta}U^{\alpha}U^{\delta}\neq 0$；对类光测地线，普适条件是说：每条以 $k^{\mu}$ 为切矢量的测地线，都必须至少有一个点使得 $k_{[\alpha}R_{\beta]\gamma\delta[\epsilon}k_{\xi]}k^{\gamma}k^{\delta}\neq 0$。这些花哨的条件不过是用来排除那些曲率在某些方向上一致为零的特殊度规。

奇点定理有许多种形式，各自从不同的一组假设出发。这个故事的寓意似乎是：广义相对论中典型的含时解通常终于奇点。（或者始于奇点；某些定理暗示了宇宙学奇点的存在，比如大爆炸。）这对 GR 来说算是个问题，因为这个理论其实并不真正适用于奇点本身，奇点的存在因此代表着描述的不完备性。对待这个问题的传统态度，是寄希望于众人苦苦追寻的量子引力理论，能以某种方式化解经典 GR 的奇点。

与此同时，我们还可以从“奇点藏在事件视界背后”这一想法中获得安慰。这一信念概括在\textbf{宇宙监督假设}（cosmic censorship conjecture）之中：

在满足主能量条件的渐近平直时空中，从通有的、初始非奇异的态出发的引力坍缩，不能形成裸奇点。

\textbf{裸奇点}（naked singularity）是指信号可以从那里到达 $\mathscr{I}^{+}$ 的奇点；也就是说，是没有藏在事件视界背后的奇点。注意，这个假设讲的是裸奇点的\emph{形成}，而不是它们的\emph{存在}；当然有一些解，其类空裸奇点存在于过去（比如 Schwarzschild 白洞），或者类时奇点存在于所有时间（比如下面要讨论的超极端带电黑洞）。宇宙监督假设至今未被证明，尽管人们已经花了大量力气去寻找有说服力的反例，却一无所获。要求初始数据在某种意义上“通有”这一点很重要，因为数值实验表明，经过精细调节的初始条件能够产生裸奇点。对某种形式的宇宙监督假设给出严格证明，仍然是经典广义相对论悬而未决的突出问题之一。\footnote{关于宇宙监督假设的评述见 R.M. Wald, “Gravitational Collapse and Cosmic Censorship,” http://arxiv.org/abs/gr-qc/9710068.}

宇宙监督（或某些等价假设）的一个推论是：经典黑洞从不缩小，只会越长越大。黑洞的大小由事件视界的面积来量度，这里指的是事件视界与一个类空切片之交的空间面积。于是我们有了 Hawking 的\textbf{面积定理}（area theorem）：

假设弱能量条件与宇宙监督成立，则渐近平直时空中未来事件视界的面积不减。

% ---- 书页 244 / p259 ----
对 Schwarzschild 黑洞，面积单调地依赖于质量，因此这个定理意味着 Schwarzschild 黑洞的质量只能增加。但对转动的黑洞就不再如此了；面积依赖于质量和角动量的某种组合，而且下面我们会看到，通过减小黑洞的自旋，我们真的可以从中提取能量。我们也可以通过量子力学的霍金辐射来减小黑洞的质量；这可以追溯到弯曲时空中的量子场论可以违反弱能量条件这一事实。

\section{Killing 视界}

在 Schwarzschild 度规中，Killing 矢量 $K=\partial_t$ 在事件视界处由类时变为类空。一般地，如果某个 Killing 矢量场 $\chi^{\mu}$ 沿某个类光超曲面 $\Sigma$ 处处类光，我们就说 $\Sigma$ 是 $\chi^{\mu}$ 的一个\textbf{Killing 视界}（Killing horizon）。注意矢量场 $\chi^{\mu}$ 必然与 $\Sigma$ 正交，因为一个类光曲面不可能有两个线性无关的类光切矢量。

Killing 视界的概念在逻辑上独立于事件视界的概念，但在具有时间平移对称性的时空中，两者紧密相关。在某些合理的条件下（下面会明确给出），我们有如下分类：

稳态、渐近平直时空中的每一个事件视界 $\Sigma$，都是某个 Killing 矢量场 $\chi^{\mu}$ 的 Killing 视界。

如果时空是静态的，$\chi^{\mu}$ 就是代表无穷远处时间平移的 Killing 矢量场 $K^{\mu}=(\partial_t)^{\mu}$。

如果时空是稳态但非静态的，它将是轴对称的，并带有转动 Killing 矢量场 $R^{\mu}=(\partial_\phi)^{\mu}$，而 $\chi^{\mu}$ 将是线性组合 $K^{\mu}+\Omega_{\mathrm{H}}R^{\mu}$，其中 $\Omega_{\mathrm{H}}$ 为某个常数。

例如，下面我们将考察转动黑洞的 Kerr 度规，其中事件视界是转动与时间平移的 Killing 矢量的某个线性组合的 Killing 视界。在 Kerr 情形，$\partial_t$ 变为类光的那个超曲面实际上是类时的，因此并不是 Killing 视界。

让我们把这套分类方案究竟在什么条件下成立说得更精确些。\footnote{相关讨论见 R. M. Wald, “The thermodynamics of black holes,” \emph{Living Rev.\ Rel.} \textbf{4}, 6 (2001), http://arxiv.org/gr-qc/9912119.}Carter 已经证明，对静态黑洞，事件视界是 $K^{\mu}$ 的 Killing 视界；这是一个纯几何的事实，即使不引用 Einstein 方程也成立。在稳态情形，如果我们假设存在一个转动 Killing 场 $R^{\mu}$，它具有这样的性质：由 $K^{\mu}$ 和 $R^{\mu}$ 张成的二维平面与一族二维曲面正交，那么事件视界将是这两个 Killing 场的某个线性组合的 Killing 视界——这同样是纯几何的论证。反过来，如果我们只假设黑洞是稳态的，那么一般说来我们无法证明事件视界
% ---- 书页 245 / p260 ----
是轴对称的。给定 Einstein 方程加上对物质场的某些条件，Hawking 得以证明：任何稳态黑洞的事件视界必定是某个矢量场的 Killing 视界，而且这类视界必定要么稳态、要么轴对称。在本章剩下的部分里，我们将当作上述分类成立那样来叙述；然而，对物质场作假设是出了名的棘手，我们应当记住，原则上有可能找到既非静态也非轴对称的黑洞，对它们来说事件视界可能就不是 Killing 视界。

需要着重指出的一点是：虽然稳态渐近平直时空的事件视界通常是 Killing 视界，但也很容易找到与事件视界毫不相干的 Killing 视界。考虑用惯性坐标表示的 Minkowski 空间，$\mathrm{d}s^{2}=-\mathrm{d}t^{2}+\mathrm{d}x^{2}+\mathrm{d}y^{2}+\mathrm{d}z^{2}$；显然这个时空中没有任何事件视界。生成 $x$ 方向推动（boost）的 Killing 矢量为
\begin{equation*}
\chi = x\partial_t + t\partial_x,
\tag{6.5}
\end{equation*}
其模长为
\begin{equation*}
\chi_\mu\chi^\mu = -x^{2} + t^{2}.
\tag{6.6}
\end{equation*}
它在类光曲面
\begin{equation*}
x = \pm t,
\tag{6.7}
\end{equation*}
上变为类光，因此这些曲面就是 Killing 视界。把推动 Killing 矢量与平移和转动 Killing 矢量组合起来，我们就能让这些视界在流形中到处移动；Killing 视界俯拾皆是。当然，在更有意思的时空中，Killing 矢量场要少一些，相应的视界（如果有的话）也会有更大的物理意义。

对每一个 Killing 视界，我们都可以关联一个叫做\textbf{表面引力}（surface gravity）的量。考虑一个带有 Killing 视界 $\Sigma$ 的 Killing 矢量 $\chi^{\mu}$。因为 $\chi^{\mu}$ 是 $\Sigma$ 的法矢量，沿着 Killing 视界它满足测地线方程，
\begin{equation*}
\chi^\mu\nabla_\mu\chi^\nu = -\kappa\chi^\nu,
\tag{6.8}
\end{equation*}
等号右边的出现是因为 $\chi^{\mu}$ 的积分曲线未必是仿射参量化的。参量 $\kappa$ 就是表面引力；它在视界上为常数，除非在 Killing 矢量为零、$\kappa$ 可以变号的“分岔二维球面”（bifurcation two-sphere）处。（例如，在 Schwarzschild 解的 Kruskal 图中心就会发生这种情况。）利用 Killing 方程 $\nabla_{(\mu}\chi_{\nu)}=0$ 以及 $\chi_{[\mu}\nabla_\nu\chi_{\sigma]}=0$ 这一事实（因为 $\chi^{\mu}$ 与 $\Sigma$ 正交），可以直截了当地导出一个漂亮的表面引力公式：
\begin{equation*}
\kappa^{2} = -\frac{1}{2}(\nabla_\mu\chi_\nu)(\nabla^\mu\chi^\nu).
\tag{6.9}
\end{equation*}
右边的表达式要在视界 $\Sigma$ 上取值。鼓励你自己动手验证一下这个公式。

% ---- 书页 246 / p261 ----
与 Killing 视界相联系的表面引力原则上带有任意性，因为我们总可以用一个实常数去缩放某个 Killing 场而得到另一个 Killing 场。在静态、渐近平直的时空中，时间平移 Killing 矢量 $K=\partial_t$ 可以通过令
\begin{equation*}
K_\mu K^\mu(r\to\infty) = -1
\tag{6.10}
\end{equation*}
来归一化。这又随之确定了任何相关联的 Killing 视界的表面引力。如果我们处在稳态时空中，那里的 Killing 视界与时间平移和转动的某个线性组合相联系，那么固定 $K=\partial_t$ 的归一化也就固定了这个线性组合，因此表面引力仍然是唯一的。

$\kappa$ 为什么叫做“表面引力”，只有当时空是静态的时候才变得清楚。在这种情形下我们有如下解释：

在静态、渐近平直的时空中，表面引力就是视界附近一位静态观者的加速度，由无穷远处的一位静态观者测得。

为了弄清这句话的含义，我们先来考虑静态观者。所谓静态观者，指的是其四速度 $U^{\mu}$ 与时间平移 Killing 场 $K^{\mu}$ 成正比的观者，
\begin{equation*}
K^\mu = V(x)U^\mu.
\tag{6.11}
\end{equation*}
由于四速度归一化为 $U_\mu U^{\mu}=-1$，函数 $V$ 不过就是 Killing 场的模长，
\begin{equation*}
V = \sqrt{-K_\mu K^\mu},
\tag{6.12}
\end{equation*}
因而从 Killing 视界处的零一直取到无穷远处的 1。$V$ 有时被称为“红移因子”（redshift factor），因为它把静态观者所测得的一个光子的发射频率与观测频率联系起来。回忆一下，四动量为 $p^{\mu}$ 的光子的守恒能量是 $E=-p_\mu K^{\mu}$，而由四速度为 $U^{\mu}$ 的观者测得的频率将是 $\omega=-p_\mu U^{\mu}$。因此
\begin{equation*}
\omega = \frac{E}{V},
\tag{6.13}
\end{equation*}
而静态观者 1 发射的一个光子将被静态观者 2 观测到具有波长 $\lambda=2\pi/\omega$，
\begin{equation*}
\lambda_2 = \frac{V_2}{V_1}\lambda_1.
\tag{6.14}
\end{equation*}
特别地，在 $V=1$ 的无穷远处，我们将观测到波长 $\lambda_{\infty}=\lambda_{1}/V_{1}$。

现在我们转向“从无穷远处看到的加速度”这个概念。静态观者一般不会沿测地线运动；比如，粒子往往要落进黑洞，而不是悬停在黑洞旁边、保持空间坐标固定。

% ---- 书页 247 / p262 ----
我们可以把四加速度 $a^{\mu}=U^{\sigma}\nabla_{\sigma}U^{\mu}$ 用红移因子表示为
\begin{equation*}
a_\mu = \nabla_\mu\ln V,
\tag{6.15}
\end{equation*}
这一点你很容易验证。加速度的大小
\begin{equation*}
a = \sqrt{a_\mu a^\mu} = V^{-1}\sqrt{\nabla_\mu V\nabla^\mu V},
\tag{6.16}
\end{equation*}
在 Killing 视界处将趋于无穷——要让一个物体保持在静态轨迹上，需要无穷大的加速度。但无穷远处的观者探测到的加速度将被因子 $V$“红移”；这个结果正是表面引力。于是我们断言，
\begin{equation*}
\kappa = Va = \sqrt{\nabla_\mu V\nabla^\mu V},
\tag{6.17}
\end{equation*}
在视界 $\Sigma$ 上取值。你可以验证这个表达式与式 (6.9) 一致。表面引力是零（$V$）与无穷（$a$）的乘积，但通常取有限值。当我们说观测到的加速度被红移了，我们脑海中的图像是：把一根检验细绳从一个位于视界处的静态物体一直拉伸到无穷远处的观者，然后测量细绳在无穷远那一端的加速度。（值得花点时间看看，你能否把这个挥手带过的论证提升为某种更严格的东西。）

如果时空是稳态但非静态的，上面的考虑会在哪里出问题？我们仍然有一个渐近的时间平移 Killing 矢量 $K=\partial_t$，而且可以像式 (6.11) 那样，把四速度与 $K^{\mu}$ 平行的观者定义为稳态观者；红移仍由式 (6.14) 给出。问题在于，$K^{\mu}$ 不会在 Killing 视界处变为类光，而一般是在视界之外某个类时曲面上变为类光。这个 $K^{\mu}K_\mu=0$ 的地方被称为\textbf{静止极限面}（stationary limit surface）（有时也称“能层面”（ergosurface）），因为在这个曲面之内 $K^{\mu}$ 是类空的，因而任何观者都无法保持静止，哪怕这里还在事件视界之外。这样的观者不得不相对于 Killing 场运动，但不必朝着黑洞的方向运动。由式 (6.12) 和式 (6.14)，当我们趋近静止极限面时，稳态观者的红移发散，因此这个曲面也被称为\textbf{无限红移面}（infinite redshift surface）。正如我们在讨论 Kerr 度规时将会看到的，静止极限面与事件视界之间的区域——能层（ergosphere）——是一个类时路径不可避免地被拖着随黑洞一起转动的地方。在稳态时空中我们将继续使用“表面引力”这个标签，并用 Killing 矢量 $\chi^{\mu}$（它在事件视界上确实变为类光）来计算它，即使所得的量已不能解释为无穷远处所看到的静态观者的引力加速度。

让我们把这些概念应用到 Schwarzschild 情形，看看它们如何运作。对于度规
\begin{equation*}
\mathrm{d}s^{2} = -\left(1-\frac{2GM}{r}\right)\mathrm{d}t^{2} + \left(1-\frac{2GM}{r}\right)^{-1}\mathrm{d}r^{2} + r^{2}\,\mathrm{d}\Omega^{2},
\tag{6.18}
\end{equation*}

% ---- 书页 248 / p263 ----
Killing 矢量和静态四速度为
\begin{equation*}
K^\mu = (1,\ 0,\ 0,\ 0), \qquad U^\mu = \left[\left(1-\frac{2GM}{r}\right)^{-1/2},\ 0,\ 0,\ 0\right],
\tag{6.19}
\end{equation*}
于是红移因子为
\begin{equation*}
V = \sqrt{1-\frac{2GM}{r}}.
\tag{6.20}
\end{equation*}
（注意这与上一章中我们对红移的计算相符。）由式 (6.15)，加速度为
\begin{equation*}
a_\mu = \frac{GM}{r^{2}\left(1-\dfrac{2GM}{r}\right)}\nabla_\mu r,
\tag{6.21}
\end{equation*}
其中当然有 $\nabla_\mu r=\delta^{r}_{\mu}$。于是加速度的大小为
\begin{equation*}
a = \frac{GM}{r^{2}\left(1-\dfrac{2GM}{r}\right)^{1/2}}.
\tag{6.22}
\end{equation*}
表面引力是 $\kappa=Va$ 在事件视界 $r=2GM$ 处的取值，而
\begin{equation*}
Va = \frac{GM}{r^{2}},
\tag{6.23}
\end{equation*}
所以 Schwarzschild 黑洞的表面引力为
\begin{equation*}
\kappa = \frac{1}{4GM}.
\tag{6.24}
\end{equation*}
表面引力随质量增大反而减小，这看上去也许令人诧异，但只要看一眼式 (6.23) 就明白是怎么回事了：在固定半径处，增大 $M$ 确实会使组合 $Va$ 变大，但质量的增大同时也使 Schwarzschild 半径增大，后一种效应占了上风。因此，大黑洞的表面引力实际上比小黑洞的还要弱；这与对潮汐力的考察相一致——黑洞越大，潮汐力也越小。

\section{质量、电荷与自旋}

我们在上面曾宣称，广义相对论最一般的稳态黑洞解由质量、电荷和自旋刻画，那么就应该考虑这些量在 GR 中如何定义。电荷最容易考虑，所以我们从它开始；更多细节见附录 E 中关于 Stokes 定理的讨论。我们将具体考察电荷，不过磁荷也可以用同样的方式来处理。

% ---- 书页 249 / p264 ----
麦克斯韦方程组把电磁场强张量 $F_{\mu\nu}$ 与电流四矢量 $J^{\mu}_{e}$ 联系起来，
\begin{equation*}
\nabla_\nu F^{\mu\nu} = J^{\mu}_{e}.
\tag{6.25}
\end{equation*}
穿过类空超曲面 $\Sigma$ 的电荷由该超曲面上坐标 $x^{i}$ 的积分给出，
\begin{align*}
Q &= -\int_{\Sigma} d^{3}x\,\sqrt{\gamma}\, n_\mu J^{\mu}_{e}\\
  &= -\int_{\Sigma} d^{3}x\,\sqrt{\gamma}\, n_\mu\nabla_\nu F^{\mu\nu},
\tag{6.26}
\end{align*}
其中 $\gamma_{ij}$ 是与 $\Sigma$ 相联系的诱导度规，$n^{\mu}$ 是单位法矢量。这个负号保证正的电荷密度加上指向未来的法矢量会给出正的总电荷。于是可以用 Stokes 定理把电荷表示成一个边界积分，
\begin{equation*}
Q = -\int_{\partial\Sigma} d^{2}x\,\sqrt{\gamma^{(2)}}\, n_\mu\sigma_\nu F^{\mu\nu},
\tag{6.27}
\end{equation*}
其中边界 $\partial\Sigma$——典型地是空间无穷远处的一个二维球面——具有度规 $\gamma^{(2)}_{ij}$ 和指向外部的法矢量 $\sigma^{\mu}$。磁荷则可以通过把 $F^{\mu\nu}$ 换成对偶张量 ${*F}^{\mu\nu}=\frac{1}{2}\epsilon^{\mu\nu\rho\sigma}F_{\rho\sigma}$ 来确定。因此，要计算总电荷，我们只需要知道电磁场在空间无穷远处的行为。在附录 E 中，我们对 Minkowski 空间中的一个点电荷作了显式计算，得到的结果不出所料，但很好地检验了我们的约定确实自洽。

现在我们转向渐近平直时空的总能量（或质量）的概念。这是一个比电荷棘手得多的概念；一方面，在广义相对论中能量-动量是张量而不是矢量；另一方面，能量-动量张量 $T_{\mu\nu}$ 只描述物质的性质，而不描述引力场的性质。但回忆一下，第 3 章中我们讨论过：如果时空是稳态的、具有类时 Killing 矢量场 $K^{\mu}$，我们仍然可以定义一个守恒的总能量。我们先构造一个流
\begin{equation*}
J^{\mu}_{T} = K_\nu T^{\mu\nu},
\tag{6.28}
\end{equation*}
其中 $T^{\mu\nu}$ 是能量-动量张量。因为这个流是无散的（由 Killing 方程和 $T^{\mu\nu}$ 的守恒而来），我们可以在一个类空曲面 $\Sigma$ 上积分，得到一个守恒的能量，
\begin{equation*}
E_{T} = \int_{\Sigma} d^{3}x\,\sqrt{\gamma}\, n_\mu J^{\mu}_{T},
\tag{6.29}
\end{equation*}
就像电荷那样。这个表达式虽然有趣，但显然有几点不足。例如考虑 Schwarzschild 度规，它有一个
% ---- 书页 250 / p265 ----
Killing 矢量，但 $T^{\mu\nu}$ 处处为零。那么 Schwarzschild 黑洞的能量就是零吗？无论从物理还是数学的理由看，都有理由怀疑并非如此；毕竟那里有一个奇点，它使得这个积分很难算出来。此外，Schwarzschild 黑洞可以由一颗具有确定非零能量的大质量恒星演化而来，我们或许愿意让这个能量守恒。因此值得去寻找另一种能量定义，以便更好地捕捉我们对黑洞时空的直观图像。

我们暂时仍然限于具有类时 Killing 矢量 $K^{\mu}$ 的时空，考虑一个新的流
\begin{equation*}
J^{\mu}_{R} = K_\nu R^{\mu\nu}.
\tag{6.30}
\end{equation*}
利用 Einstein 方程，我们也可以把它等价地写成
\begin{equation*}
J^{\mu}_{R} = 8\pi G K_\nu\left(T^{\mu\nu}-\frac{1}{2}T g^{\mu\nu}\right).
\tag{6.31}
\end{equation*}
Ricci 张量并不是无散的；我们有的是缩并的 Bianchi 恒等式，
\begin{equation*}
\nabla_\mu R^{\mu\nu} = \frac{1}{2}\nabla^{\nu}R.
\tag{6.32}
\end{equation*}
但这个恒等式加上 Killing 方程，就足以保证我们的新流是守恒的。为看清这一点，我们只需计算
\begin{equation*}
\nabla_\mu J^{\mu}_{R} = (\nabla_\mu K_\nu)R^{\mu\nu} + K_\nu(\nabla_\mu R^{\mu\nu}).
\tag{6.33}
\end{equation*}
第一项为零，因为 $R^{\mu\nu}$ 是对称的，而 $\nabla_\mu K_\nu$ 是反对称的（由 Killing 方程）。于是利用式 (6.32)，我们得到
\begin{equation*}
\nabla_\mu J^{\mu}_{R} = \frac{1}{2}K_\nu\nabla^{\nu}R = 0,
\tag{6.34}
\end{equation*}
我们知道它为零，因为 $R$ 沿 Killing 矢量的方向导数为零，见式 (3.178)。

和前面一样，我们可以定义一个与这个流相联系的守恒能量，
\begin{equation*}
E_{R} = \frac{1}{4\pi G}\int_{\Sigma} d^{3}x\,\sqrt{\gamma}\, n_\mu J^{\mu}_{R},
\tag{6.35}
\end{equation*}
其中归一化的选取是为了将来的方便。能量 $E_{R}$ 将不依赖于类空超曲面 $\Sigma$ 的选取，因而是守恒的。这个能量概念比起 $E_{T}$ 有一项显著的优势，源于这样一个事实：$E_{R}$ 可以改写成空间无穷远处一个二维球面上的面积分。为看清这一点，回忆式 (3.177)：任何 Killing 矢量都满足 $\nabla_\mu\nabla_\nu K^{\mu}=K^{\mu}R_{\mu\nu}$；于是这个流本身可以写成一个全导数，
\begin{equation*}
J^{\mu}_{R} = \nabla_\nu(\nabla^{\mu}K^{\nu}),
\tag{6.36}
\end{equation*}
 于是有
% ---- 书页 251 / p266 ----
\begin{equation*}
E_{R} = \frac{1}{4\pi G}\int_{\Sigma} d^{3}x\,\sqrt{\gamma}\, n_{\mu}\nabla_{\nu}\left(\nabla^{\mu}K^{\nu}\right).
\tag{6.37}
\end{equation*}

注意，由 Killing 方程升起指标可得 $\nabla^{\mu}K^{\nu}=-\nabla^{\nu}K^{\mu}$。因此，正如前面处理电荷时那样，我们可以再次利用 Stokes 定理，把 $E_{R}$ 写成空间无穷远处的积分，
\begin{equation*}
\boxed{\,E_{R} = \frac{1}{4\pi G}\int_{\partial\Sigma} d^{2}x\,\sqrt{\gamma^{(2)}}\, n_{\mu}\sigma_{\nu}\nabla^{\mu}K^{\nu}.\,}
\tag{6.38}
\end{equation*}
这个表达式就是与类时 Killing 矢量 $K^{\mu}$ 相联系的\textbf{Komar 积分}（Komar integral）；它可以解释为稳态时空的总能量。

为了让我们自己确信方向没错，就用度规为式 (6.18) 的 Schwarzschild 时空来算一算 Komar 积分。归一化为 $n_{\mu}n^{\mu}=-1$ 与 $\sigma_{\mu}\sigma^{\mu}=+1$ 的法矢量，其非零分量为
\begin{equation*}
n_{0} = -\left(1-\frac{2GM}{r}\right)^{1/2}, \qquad \sigma_{1} = \left(1-\frac{2GM}{r}\right)^{-1/2},
\tag{6.39}
\end{equation*}
其余分量均为零。于是我们有
\begin{equation*}
n_{\mu}\sigma_{\nu}\nabla^{\mu}K^{\nu} = -\nabla^{0}K^{1}.
\tag{6.40}
\end{equation*}
Killing 矢量为 $K^{\mu}=(1,0,0,0)$，因此我们可以毫不费力地算出
\begin{align*}
\nabla^{0}K^{1} &= g^{00}\nabla_{0}K^{1}\\
&= g^{00}\left(\partial_{0}K^{1}+\Gamma^{1}_{0\lambda}K^{\lambda}\right)\\
&= g^{00}\Gamma^{1}_{00}K^{0}\\
&= -\left(1-\frac{2GM}{r}\right)^{-1}\frac{GM}{r^{2}}\left(1-\frac{2GM}{r}\right)\\
&= -\frac{GM}{r^{2}}.
\tag{6.41}
\end{align*}
无穷远处二维球面上的度规是
\begin{equation*}
\gamma^{(2)}_{ij}\,\mathrm{d}x^{i}\mathrm{d}x^{j} = r^{2}\left(\mathrm{d}\theta^{2}+\sin^{2}\theta\,\mathrm{d}\phi^{2}\right),
\tag{6.42}
\end{equation*}
于是
\begin{equation*}
\sqrt{\gamma^{(2)}} = r^{2}\sin\theta.
\tag{6.43}
\end{equation*}

% ---- 书页 252 / p267 ----
把这一切合在一起，Schwarzschild 黑洞的能量就是
\begin{align*}
E_{R} &= \frac{1}{4\pi G}\int d\theta\, d\phi\, r^{2}\sin\theta\left(\frac{GM}{r^{2}}\right)\\
&= M.
\tag{6.44}
\end{align*}
这当然正是我们想要的结果，同时也解释了式 (6.35) 中所选取的归一化。

尽管得到了正确的答案，我们还是应该想一想刚才到底发生了什么。具体来说，我们是把流 $J^{\mu}_{R}=K_{\nu}R^{\mu\nu}$ 在一个类空切片上积分而得到这个能量的，并且发现结果可以写成空间无穷远处的积分。但对 Schwarzschild 来说，度规满足真空 Einstein 方程 $R_{\mu\nu}=0$；这样一来，要从对 $J^{\mu}_{R}$ 的积分中得到非零的结果似乎就很困难，正如式 (6.29) 那样。可是如果想一想最大延拓 Schwarzschild 解的结构，我们就会意识到，可以画出两类类空切片：一类穿过虫洞一直延伸到第二个渐近区域，另一类则终止于奇点。如果切片穿过虫洞，另一个渐近区域就会为 $\partial\Sigma$ 提供另一个组成部分，从而对式 (6.38) 提供另一份贡献；这份贡献会恰好抵消，于是总能量确实会是零。如果切片撞上了奇点，我们就不知道该怎么处理了。尽管如此，无论哪种情形，把我们的结果 (6.44) 当作正确的答案都是合理的。关键在于，式 (6.38) 只涉及空间无穷远处的贡献，所以无论内部发生什么，它都应当是能量的一个有效表达式。我们甚至可以设想内部出现随时间变化的行为；只要 $K^{\mu}$ \emph{渐近地}是类时 Killing 矢量，Komar 能量就是良定义的。比如说，我们可以考虑从一颗初始静态的恒星出发的球对称引力坍缩。直接在 $\Sigma$ 上计算积分 (6.35)，会给出一个合理的总质量答案，而且当恒星坍缩为黑洞时，这个总质量不应改变（我们设想的是球对称，因此引力辐射无法把能量带到无穷远）。于是，在坍缩之前有效的 Komar 积分 (6.38)，在坍缩成黑洞之后也尽可以放心地解释为能量。当然，出于某些目的，我们或许想容许通过引力辐射损失能量，那种情形下就必须小心对待如何把切片延伸到无穷远；人们可以在未来类光无限远处定义一个“Bondi 质量”（Bondi mass），用它来追踪通过辐射损失的能量。

关于 Komar 公式的另一个疑虑是，它是否真的是我们应当认作“能量”的那种东西——能量通常是指与时间平移不变性相联系的守恒量。支持这种解释的最好理由很简单：$E_{R}$ 无疑是某种守恒量，而且它与我们认为应当是 Schwarzschild 能量的数值相一致（在 Newton 极限下对一团质量合集也是如此，你可以自行验证），那么它还能是什么呢？另一种做法是考虑广义相对论的哈密顿表述，在渐近平直时空中仔细定义时间平移的生成元，再把它认同为总能量。这项工作最先由 Arnowitt、Deser 和 Misner 完成，他们的结果就是所谓的\textbf{ADM 能量}（ADM energy）。在一个
% ---- 书页 253 / p268 ----
渐近平直时空中，我们可以像在微扰论中那样把度规写成
\begin{equation*}
g_{\mu\nu} = \eta_{\mu\nu} + h_{\mu\nu},
\tag{6.45}
\end{equation*}
只不过这里我们只要求分量 $h_{\mu\nu}$ 在空间无穷远处很小，而不必处处都小。于是 ADM 能量可以写成空间无穷远处一个二维球面上的积分：
\begin{equation*}
E_{\mathrm{ADM}} = \frac{1}{16\pi G}\int_{\partial\Sigma} d^{2}x\,\sqrt{\gamma^{(2)}}\,\sigma^{i}\left(\partial_{j}h^{j}{}_{i}-\partial_{i}h^{j}{}_{j}\right),
\tag{6.46}
\end{equation*}
其中空间指标用 $\delta^{ij}$（无穷远处的空间度规）升起。这个公式看上去依赖于坐标的选取，但在我们的假设之下实际上是良定义的。如果 $h_{\mu\nu}$ 在无穷远处不依赖于时间，可以验证 ADM 能量与 Komar 能量实际上是一致的。这让我们更加相信 Komar 积分确实表示能量。不过，在某种意义上 ADM 能量更为“体面”；例如，如果存在与引力非最小耦合的长程标量场，Komar 积分就可能遇到麻烦。但就我们眼前的目的而言，Komar 能量已经完全够用了。

我们希望某个被叫作“能量”的东西具有这样一种品质：对任何物理位形它都是正的；否则，一个零能量的态就可以衰变成一些正能量的碎片和一些负能量的碎片。第 4 章讨论过的能量条件为物质场给出了正能量的概念，但我们可能会担心，引力的负能量贡献会带来麻烦。所幸在 GR 中我们有\textbf{正能量定理}（positive energy theorem），它最先由 Schoen 和 Yau（译注：原文排印作 “Shoen”，应为 Schoen。）证明：

无奇点的、满足 Einstein 方程和主能量条件的渐近平直时空，其 ADM 能量是非负的。而且，Minkowski 时空是唯一使 ADM 能量为零的此类时空。

如果容许奇点存在，反例显然是有的，比如 $M<0$ 的 Schwarzschild 解。不过，如果一个带奇点的时空（比如 $M>0$ 的 Schwarzschild 解）是作为无奇点初数据的演化结果而达到的，定理就仍然适用。这样看来，在广义相对论中我们似乎是安全的，不会遇到负能量的孤立系统。

最后，我们可以转向自旋（角动量）；讨论完能量之后，这一步轻而易举。设想我们有一个转动 Killing 矢量 $R=\partial_{\phi}$。与时间平移的情形完全类似，我们可以定义一个守恒流
\begin{equation*}
J^{\mu}_{\phi} = R_{\nu}R^{\mu\nu},
\tag{6.47}
\end{equation*}
由它将得到把守恒角动量 $J$ 表示成空间无穷远处积分的表达式，

% ---- 书页 254 / p269 ----
\begin{equation*}
\boxed{\,J = -\frac{1}{8\pi G}\int_{\partial\Sigma} d^{2}x\,\sqrt{\gamma^{(2)}}\, n_{\mu}\sigma_{\nu}\nabla^{\mu}R^{\nu}.\,}
\tag{6.48}
\end{equation*}
（“J”既被用来表示流又被用来表示角动量，“R”既是转动 Killing 矢量又是 Ricci 张量，这都挺遗憾。可是字母就只有这么多，只能将就了。）与能量一样，即使 $R^{\mu}$ 只是渐近地才是 Killing 矢量，这个表达式仍然有效。注意，这里的归一化与能量积分中的不同；例如，可以通过对弱引力场中缓慢运动的物体求出该表达式的值，来为这个归一化提供依据。

\section{带电（Reissner–Nordström）黑洞}

现在我们转向表示带电黑洞的精确解。这类解与现实的天体物理情形关系并不太大；在真实世界中，一个高电荷黑洞会通过与黑洞附近物质的相互作用迅速被中和。不过，带电黑洞仍然展示了更一般情形的若干重要特征。在这个问题中，完整的球对称性依然存在；因此我们知道可以把度规写成
\begin{equation*}
\mathrm{d}s^{2} = -e^{2\alpha(r,t)}\mathrm{d}t^{2} + e^{2\beta(r,t)}\mathrm{d}r^{2} + r^{2}\mathrm{d}\Omega^{2}.
\tag{6.49}
\end{equation*}
然而现在我们不再处于真空，因为黑洞会带有非零的电磁场，而电磁场又充当能量-动量的源。电磁场的能量-动量张量为
\begin{equation*}
T_{\mu\nu} = F_{\mu\rho}F_{\nu}{}^{\rho} - \frac{1}{4}g_{\mu\nu}F_{\rho\sigma}F^{\rho\sigma},
\tag{6.50}
\end{equation*}
其中 $F_{\mu\nu}$ 是电磁场强张量。既然问题具有球对称性，最一般的场强张量就只能有如下分量
\begin{gather*}
F_{tr} = f(r,t) = -F_{rt}\\
F_{\theta\phi} = g(r,t)\sin\theta = -F_{\phi\theta},
\tag{6.51}
\end{gather*}
其中 $f(r,t)$ 和 $g(r,t)$ 是两个有待场方程确定的函数，没有写出的分量均为零。$F_{tr}$ 对应于径向电场，而 $F_{\theta\phi}$ 对应于径向磁场。如果你正纳闷那个 $\sin\theta$ 是从哪儿来的，回忆一下：不应当依赖于 $\theta$ 和 $\phi$ 的量，是磁场的径向分量 $B^{r}=\epsilon^{01\mu\nu}F_{\mu\nu}$。对球对称的度规，
\begin{equation*}
\epsilon^{\rho\sigma\mu\nu} = \frac{1}{\sqrt{-g}}\tilde{\epsilon}^{\rho\sigma\mu\nu}
\end{equation*}
与 $(\sin\theta)^{-1}$ 成正比，所以我们希望在 $F_{\theta\phi}$ 中带有一个 $\sin\theta$ 因子。

% ---- 书页 255 / p270 ----
这种情形下的场方程既有 Einstein 方程，也有麦克斯韦方程组：
\begin{gather*}
g^{\mu\nu}\nabla_{\mu}F_{\nu\sigma} = 0\\
\nabla_{[\mu}F_{\nu\rho]} = 0.
\tag{6.52}
\end{gather*}
这两组方程彼此耦合在一起，因为电磁场强张量通过能量-动量张量进入 Einstein 方程，而度规又显式地出现在麦克斯韦方程组之中。

然而，这些困难并非不可克服；与真空情形中我们如法炮制的步骤相类似的做法，同样会给出带电情形的解。我们不会显式地写出各个步骤，只把最后的结果抄下来。这个解就是\textbf{Reissner–Nordström（RN）度规}（Reissner–Nordström metric），
\begin{equation*}
\boxed{\,\mathrm{d}s^{2} = -\Delta\,\mathrm{d}t^{2} + \Delta^{-1}\mathrm{d}r^{2} + r^{2}\mathrm{d}\Omega^{2},\,}
\tag{6.53}
\end{equation*}
其中
\begin{equation*}
\Delta = 1 - \frac{2GM}{r} + \frac{G(Q^{2}+P^{2})}{r^{2}}.
\tag{6.54}
\end{equation*}
在这个表达式中，$M$ 依旧解释为黑洞的质量；$Q$ 是总电荷，$P$ 是总磁荷。孤立的磁荷（磁单极子）从未在自然界中被观测到，但这并不妨碍我们写下它们如果真的存在时将会产生的度规。\footnote{本章采用的单位制中，库仑定律里没有 $4\pi$ 因子。要想与其他章节比较，把每一处出现的 $Q$ 或 $P$ 除以 $4\pi$ 即可。}有很好的理论理由让我们相信，如果各种相互作用在极高能量下“大统一”，磁单极子就可能存在，但它们必定非常重，而且极为稀少。当然，即使不存在任何磁单极子，黑洞也可能带上磁荷。实际上，电荷与磁荷进入度规的方式完全相同，因此在表达式中保留 $P$ 并不会引入任何额外的复杂之处。保守的读者尽可以随意令 $P=0$。与这个解相联系的电磁场为
\begin{gather*}
E_{r} = F_{rt} = \frac{Q}{r^{2}}\\
B_{r} = \frac{F_{\theta\phi}}{r^{2}\sin\theta} = \frac{P}{r^{2}}.
\tag{6.55}
\end{gather*}
这些场的 $1/r^{2}$ 依赖关系正是我们在平直空间中所熟悉的；当然，我们知道这依赖于我们对径向坐标的精确选取。RN 度规在 $r=0$ 处有一个名副其实的曲率奇点，这可以通过计算曲率不变标量 $R_{\mu\nu\rho\sigma}R^{\mu\nu\rho\sigma}$ 来验证。视界结构
% ---- 书页 256 / p271 ----
却比 Schwarzschild 情形复杂得多。在上面关于事件视界的讨论中我们曾提出，如果我们聪明地选取了坐标，使得这个条件在 $r$ 的某个固定值处得到满足，那么 $g^{rr}=0$ 就会是定位事件视界的一个有用判据。幸运的是，式 (6.53) 的坐标恰好具有这种性质，于是事件视界将位于
\begin{equation*}
g^{rr}(r) = \Delta(r) = 1 - \frac{2GM}{r} + \frac{G(Q^{2}+P^{2})}{r^{2}} = 0.
\tag{6.56}
\end{equation*}
这将发生在
\begin{equation*}
r_{\pm} = GM \pm \sqrt{G^{2}M^{2} - G(Q^{2}+P^{2})}.
\tag{6.57}
\end{equation*}
\begin{figure}[H]
\centering
\includegraphics[width=0.72\textwidth]{fig_6.2.pdf}
\caption*{图 6.2\quad Reissner–Nordström 解的函数 $\Delta(r)=1-2GM/r+G(Q^{2}+P^{2})/r^{2}$；零点标示事件视界所在的位置。}
\end{figure}
如图 6.2 所示，随着 $GM^{2}$ 与 $Q^{2}+P^{2}$ 相对大小的不同，这可能给出两个、一个或零个解。于是我们逐个考察各种情形。

\subsection*{情形一：$GM^{2}<Q^{2}+P^{2}$}

这种情形下，系数 $\Delta$ 总为正（永不为零），度规在 $(t,r,\theta,\phi)$ 坐标中一直直到 $r=0$ 都完全正则。坐标 $t$ 总是类时的，而 $r$ 总是类空的。但 $r=0$ 处依然有奇点，只是现在它是一条类时线。由于不存在事件视界，观者前往奇点、看过一眼再返回报告所见，没有任何阻碍。这就是前面讨论过的裸奇点。对测地线的仔细分析表明，奇点是排斥的——类时测地线永远不会与 $r=0$ 相交；它们只是靠近，然后掉头离去。（类光测地线可以到达奇点，非测地线的类时曲线也可以。）

% ---- 书页 257 / p272 ----
\begin{figure}[H]
\centering
\includegraphics[width=0.72\textwidth]{fig_6.3.pdf}
\caption*{图 6.3\quad Reissner–Nordström 解在 $Q^{2}+P^{2}>GM^{2}$ 时的共形图。原点处有一个裸奇点。}
\end{figure}

当 $r\to\infty$ 时解趋近于平直时空，而且正如我们刚刚看到的，因果结构看上去处处正常。因此共形图就会与 Minkowski 空间的一模一样，只不过现在 $r=0$ 是一个奇点，如图 6.3 所示。

奇点的裸露冒犯了我们的体面感，也同样冒犯了宇宙监督假设。实际上，我们决不应指望引力坍缩的结果会是一个 $GM^{2}<Q^{2}+P^{2}$ 的黑洞。粗略地说，这个条件说的是，黑洞的总能量比单独电磁场贡献的能量还要少——也就是说，携带电荷的那些物质，其质量必须为负。因此这个解通常被认为是非物理的。还要注意，这个时空中不存在任何 Cauchy 面，因为类时线可以在奇点处开始或终止。

\subsection*{情形二：$GM^{2}>Q^{2}+P^{2}$}

我们预期这种情形适用于现实的引力坍缩；电磁场中的能量小于总能量。此时度规系数 $\Delta(r)$ 在 $r$ 很大和很小时都为正，而在两个零点 $r_{\pm}=GM\pm\sqrt{G^{2}M^{2}-G(Q^{2}+P^{2})}$ 之间为负。度规在 $r_{+}$ 和 $r_{-}$ 两处都有坐标奇点；与 Schwarzschild 的情形一样，这两处都可以通过坐标变换消去。

由 $r=r_{\pm}$ 定义的曲面都是类光的，而且二者都是事件视界。$r=0$ 处的奇点是一条类时线，而不像 Schwarzschild 中那样是一个类空曲面。如果你是一个从远处落入黑洞的观者，$r_{+}$ 就好比 Schwarzschild 度规中的 $2GM$：在这个半径处，$r$ 从类空坐标切换成类时坐标，你必然朝 $r$ 减小的方向运动。黑洞外面的目击者也会看到同样的
% ---- 书页 258 / p273 ----
现象，与不带电黑洞外面的情形相同——落入的观者看上去运动得越来越慢，红移也越来越厉害。

但是，从 $r_{+}$ 一路坠向越来越小的半径这种不可避免的坠落，只会持续到你抵达类光曲面 $r=r_{-}$ 为止；在那里 $r$ 重新变回类空坐标，朝 $r$ 减小方向的运动便可以止住。因此你并不一定要撞上 $r=0$ 处的奇点；这是意料之中的，因为 $r=0$ 是一条类时线（因而未必处在你的未来之中）。实际上，你既可以选取继续走向 $r=0$，也可以开始朝 $r$ 增大的方向折返，重新穿过 $r=r_{-}$ 处的类光曲面。此后 $r$ 将再一次成为类时坐标，只不过取向反了过来；你被迫朝 $r$ \emph{增大}的方向运动。最终你会再一次被抛过 $r=r_{+}$，就仿佛从一个白洞里冒出来、进入宇宙的其余部分。从这里你又可以选取回到黑洞中去——这一次是另一个洞，与你最初进入的那个不同——并且随你喜欢把这样的旅行重复多少次。这个小故事对应于图 6.4 中的共形图；当然，通过选取适当的坐标并把 Reissner–Nordström 度规解析延拓到极致，可以更严格地把它推导出来。

\begin{figure}[H]
\centering
\includegraphics[width=0.72\textwidth]{fig_6.4.pdf}
\caption*{图 6.4\quad Reissner–Nordström 解在 $GM^{2}>Q^{2}+P^{2}$ 时的共形图。黑洞外部区域有无穷多个副本。}
\end{figure}

% ---- 书页 259 / p274 ----
\begin{figure}[H]
\centering
\includegraphics[width=0.72\textwidth]{fig_6.5.pdf}
\caption*{图 6.5\quad 极端 Reissner–Nordström 解（$GM^{2}=Q^{2}+P^{2}$）的共形图。原点处有一个裸奇点，外部区域则有无穷多个。}
\end{figure}

这番情景里有多少是科学、有多少是科学幻想呢？大概科学不多。如果你从一位身处黑洞内部、正要穿过 $r=r_{-}$ 处事件视界的观者的角度来考虑这个世界，你会注意到，这位观者可以回溯时间，看到外部（渐近平直）宇宙的整部历史——至少从黑洞看来是如此。但是，观者是在自己有限的固有时间内看到这段（无限长的）历史的——于是，当观者接近 $r_{-}$ 时，到达他们的任何信号都会被无限蓝移。因此，任何进入 RN 黑洞的非球对称微扰，都很可能会剧烈地扰动我们所描述的几何。真实的几何会是什么样子很难说，但没有很好的理由相信它必定包含无穷多个通过种种虫洞彼此连通的渐近平直区域。\footnote{关于这个问题的一些研究，见 E. Poisson and W. Israel, \emph{Phys.\ Rev.\ D} \textbf{41}, 1796 (1990).}

\subsection*{情形三：$GM^{2}=Q^{2}+P^{2}$}

这种情形称为\textbf{极端}（extremal）Reissner–Nordström 解。一方面，极端黑洞是一件有趣的理论玩具；在研究黑洞在量子引力中扮演的角色时，人们经常考察这个解。在超对称理论中，极端黑洞可以让某些对称性保持不破缺，这给计算带来极大的方便。另一方面，它看上去是不稳定的，因为只要加入一点点物质，它就会变成情形二。

极端黑洞的 $\Delta(r)$ 只在单个半径 $r=GM$ 处为零。这代表一个事件视界，但 $r$ 坐标永远不是类时的；它在 $r=GM$ 处变成类光的，而在两侧都是类空的。与其他情形一样，$r=0$ 处的奇点是一条类时线。所以对这个黑洞来说，你同样可以避开奇点，继续向未来运动到渐近平直区域的额外副本中去，只不过奇点总是“在左边”。共形图如图 6.5 所示。

极端黑洞的一个迷人性质是，质量在某种意义上被电荷平衡了。更具体地说，两个携带同号电荷的极端黑洞会彼此产生引力吸引，却又彼此产生电磁排斥，而这两种效应恰好抵消。事实上，对于相耦合的 Einstein–Maxwell 方程组，我们能找到\emph{精确}解，表示处在稳定位形中的任意数目的此类黑洞。为看清这一点，先回到 Reissner–Nordström 度规本身；为简单起见，我们只考虑电荷而不考虑磁荷。在极端性条件下 $GM^{2}=Q^{2}$，度规化为
\begin{equation*}
\mathrm{d}s^{2} = -\left(1-\frac{GM}{r}\right)^{2}\mathrm{d}t^{2} + \left(1-\frac{GM}{r}\right)^{-2}\mathrm{d}r^{2} + r^{2}\mathrm{d}\Omega^{2}.
\tag{6.58}
\end{equation*}
通过定义一个平移的径向坐标
\begin{equation*}
\rho = r - GM,
\tag{6.59}
\end{equation*}

% ---- 书页 260 / p275 ----
度规便取各向同性的形式
\begin{equation*}
\mathrm{d}s^{2} = -H^{-2}(\rho)\mathrm{d}t^{2} + H^{2}(\rho)\left[\mathrm{d}\rho^{2} + \rho^{2}\mathrm{d}\Omega^{2}\right],
\tag{6.60}
\end{equation*}
其中
\begin{equation*}
H(\rho) = 1 + \frac{GM}{\rho}.
\tag{6.61}
\end{equation*}
因为 $\mathrm{d}\rho^{2}+\rho^{2}\mathrm{d}\Omega^{2}$ 正是三个空间维度中的平直度规，我们同样可以把式 (6.60) 写成
\begin{equation*}
\mathrm{d}s^{2} = -H^{-2}(\vec{x})\mathrm{d}t^{2} + H^{2}(\vec{x})\left[\mathrm{d}x^{2} + \mathrm{d}y^{2} + \mathrm{d}z^{2}\right],
\tag{6.62}
\end{equation*}
其中 $H$ 可以写成
\begin{equation*}
H = 1 + \frac{GM}{|\vec{x}|}.
\tag{6.63}
\end{equation*}
在原先的 $r$ 坐标中，极端解的电场可以用矢势 $A_{\mu}$ 表示为
\begin{equation*}
E_{r} = F_{rt} = \frac{Q}{r^{2}} = \partial_{r}A_{0},
\tag{6.64}
\end{equation*}
其中矢势的类时分量为
\begin{equation*}
A_{0} = -\frac{Q}{r},
\tag{6.65}
\end{equation*}
并且我们设想空间分量为零（已经把磁场设为零）。在新的 $\rho$ 坐标下，再加上极端性条件 $Q^{2}=GM^{2}$，上式变为
\begin{equation*}
A_{0} = -\frac{\sqrt{G}M}{\rho+GM},
\tag{6.66}
\end{equation*}
或者等价地写成
\begin{equation*}
\sqrt{G}\,A_{0} = H^{-1}-1.
\tag{6.67}
\end{equation*}
但现在让我们忘掉 $H$ 满足式 (6.61) 这件事，径直把度规 (6.62) 和静电势 (6.67) 代入 Einstein 方程与麦克斯韦方程组，设想 $H$ 不依赖于时间（$\partial_{0}H=0$）但除此之外不受任何约束。我们可以直截了当地证明（见习题）：任何一个不依赖于时间、且满足
\begin{equation*}
\nabla^{2}H = 0,
\tag{6.68}
\end{equation*}
（其中 $\nabla^{2}=\partial_{x}^{2}+\partial_{y}^{2}+\partial_{z}^{2}$）的函数 $H(\vec{x})$，都能让这两组方程同时得到满足。这其实就是拉普拉斯方程，而且很容易写出所有在无穷远处表现良好的解；它们的形式为

% ---- 书页 261 / p276 ----
\begin{equation*}
H = 1 + \sum_{a=1}^{N}\frac{GM_{a}}{|\vec{x}-\vec{x}_{a}|},
\tag{6.69}
\end{equation*}
其中 $\vec{x}_{a}$ 给出某组 $N$ 个空间点。这些点描述了 $N$ 个极端 RN 黑洞所在的位置，各黑洞的质量为 $M_{a}$，电荷为 $Q_{a}=\sqrt{G}M_{a}$。这个多极端黑洞度规无疑是 Einstein 方程最引人注目的精确解之一。

\section{转动（Kerr）黑洞}

关于带电解，我们本来还可以多谈不少细节，但我们还是接着讲转动黑洞吧。在这种情形下寻找度规的精确解要困难得多，因为我们已经放弃了球对称性。我们转而寻找绕转动轴具有轴对称性、同时又是稳态的（即存在一个类时 Killing 矢量的）解。虽然 Schwarzschild 解和 Reissner–Nordström 解在广义相对论问世后不久就被发现了，但转动黑洞的解直到 1963 年才由 Kerr 找到。他的结果——\textbf{Kerr 度规}（Kerr metric）——就是下面这团乱麻：
\begin{equation*}
\boxed{\begin{aligned}
\mathrm{d}s^{2} &= -\left(1-\frac{2GMr}{\rho^{2}}\right)\mathrm{d}t^{2} - \frac{2GMar\sin^{2}\theta}{\rho^{2}}\left(\mathrm{d}t\,\mathrm{d}\phi + \mathrm{d}\phi\,\mathrm{d}t\right)\\
&\quad + \frac{\rho^{2}}{\Delta}\mathrm{d}r^{2} + \rho^{2}\mathrm{d}\theta^{2} + \frac{\sin^{2}\theta}{\rho^{2}}\left[(r^{2}+a^{2})^{2} - a^{2}\Delta\sin^{2}\theta\right]\mathrm{d}\phi^{2},
\end{aligned}}
\tag{6.70}
\end{equation*}
其中
\begin{equation*}
\boxed{\,\Delta(r) = r^{2} - 2GMr + a^{2}\,}
\tag{6.71}
\end{equation*}
以及
\begin{equation*}
\boxed{\,\rho^{2}(r,\theta) = r^{2} + a^{2}\cos^{2}\theta.\,}
\tag{6.72}
\end{equation*}
两个常数 $M$ 和 $a$ 参数化了各种可能的解。要验证质量 $M$ 确实等于 Komar 能量 (6.38) 并不难，只是繁琐；而 $a$ 是单位质量的角动量，
\begin{equation*}
a = J/M,
\tag{6.73}
\end{equation*}
其中 $J$ 是 Komar 角动量 (6.48)。要把电荷和磁荷 $Q$ 与 $P$ 也包括进来很容易，只需把 $2GMr$ 换成 $2GMr-G(Q^{2}+$

% ---- 书页 262 / p277 ----
$P^{2})$；得到的结果就是\textbf{Kerr–Newman 度规}（Kerr–Newman metric）。与之相联系的 1-形式势具有如下非零分量
\begin{equation*}
A_{t} = \frac{Qr - Pa\cos\theta}{\rho^{2}}, \qquad A_{\phi} = \frac{-Qar\sin^{2}\theta + P(r^{2}+a^{2})\cos\theta}{\rho^{2}}.
\tag{6.74}
\end{equation*}
所有本质性的现象在电荷不存在时依然保持，所以从现在起我们令 $Q=P=0$。

坐标 $(t,r,\theta,\phi)$ 就是所谓的\textbf{Boyer–Lindquist 坐标}（Boyer–Lindquist coordinates）。很容易验证，当 $a\to 0$ 时它们退化为 Schwarzschild 坐标。然而，如果保持 $a$ 不动而让 $M\to 0$，我们得到的仍是平直时空，只不过不是用普通的极坐标来表示的。此时度规变为
\begin{equation*}
\mathrm{d}s^{2} = -\mathrm{d}t^{2} + \frac{r^{2}+a^{2}\cos^{2}\theta}{r^{2}+a^{2}}\mathrm{d}r^{2} + \left(r^{2}+a^{2}\cos^{2}\theta\right)^{2}\mathrm{d}\theta^{2} + \left(r^{2}+a^{2}\right)\sin^{2}\theta\,\mathrm{d}\phi^{2},
\tag{6.75}
\end{equation*}
我们认出它的空间部分正是用椭球坐标表示的平直空间，如图 6.6 所示。这些坐标与欧几里得三维空间中的笛卡尔坐标之间的关系是
\begin{align*}
x &= \left(r^{2}+a^{2}\right)^{1/2}\sin\theta\cos\phi\\
y &= \left(r^{2}+a^{2}\right)^{1/2}\sin\theta\sin\phi\\
z &= r\cos\theta.
\tag{6.76}
\end{align*}

度规 (6.70) 有两个 Killing 矢量，二者都是显而易见的；由于度规系数不依赖于 $t$ 和 $\phi$，$K=\partial_{t}$ 与 $R=\partial_{\phi}$ 都是 Killing 矢量。当然，$R^{\mu}$ 表达的是解的轴对称性。矢量 $K^{\mu}$ 并不与 $t=$ 常数的超曲面正交，实际上也不与任何超曲面正交；因此这个度规是稳态的，但不是
% 本块末句在原书书页263顶部继续（“static. This makes sense; ...”），下一块请用 %JOIN% 续接本句。
静态的。这是有道理的：黑洞在转动，所以它不是静态的；但它在所有时刻都以完全相同的方式转动，所以它是稳态的。或者也可以这样理解：这个度规不可能是静态的，因为它不具有时间反演不变性——时间反演会把黑洞的角动量反过来。

\begin{figure}[H]
\centering
\includegraphics[width=0.72\textwidth]{fig_6.6.pdf}
\caption*{图 6.6\quad Kerr 度规中使用的椭球坐标 $(r,\theta)$。$r=0$ 是一个二维圆盘；$r=0$ 与 $\theta=\pi/2$ 的交集是这个圆盘边界上的圆环。}
\end{figure}

% ---- 书页 263 / p278 ----（本页第一段已在上面 tail 中接续，正文自第二段起）

Kerr 度规还拥有一个 Killing 张量（Killing tensor）。第 3 章中曾把它定义为任何满足如下条件的对称 $(0,n)$ 张量 $\sigma_{\mu_{1}\cdots\mu_{n}}$：
\begin{equation*}
\nabla_{(\lambda}\sigma_{\mu_{1}\cdots\mu_{n})} = 0.
\tag{6.77}
\end{equation*}
在 Kerr 几何中，我们可以定义 $(0,2)$ 张量
\begin{equation*}
\sigma_{\mu\nu} = 2\rho^{2}l_{(\mu}n_{\nu)} + r^{2}g_{\mu\nu}.
\tag{6.78}
\end{equation*}
在这个表达式中，两个矢量 $l$ 和 $n$（升指标后）为
\begin{gather*}
l^{\mu} = \frac{1}{\Delta}\left(r^{2}+a^{2},\, \Delta,\, 0,\, a\right)\\
n^{\mu} = \frac{1}{2\rho^{2}}\left(r^{2}+a^{2},\, -\Delta,\, 0,\, a\right).
\tag{6.79}
\end{gather*}
这两个矢量都是类光的，并且满足
\begin{equation*}
l^{\mu}l_{\mu} = 0,\qquad n^{\mu}n_{\mu} = 0,\qquad l^{\mu}n_{\mu} = -1.
\tag{6.80}
\end{equation*}
有了这些定义，你可以自己验证 $\sigma_{\mu\nu}$ 是一个 Killing 张量。

我们为 Kerr 选定的坐标，使得事件视界出现在那些让 $g^{rr}=0$ 的固定 $r$ 值处。由于 $g^{rr}=\Delta/\rho^{2}$，而 $\rho^{2}\geq 0$，这发生在
\begin{equation*}
\Delta(r) = r^{2} - 2GMr + a^{2} = 0.
\tag{6.81}
\end{equation*}
与 Reissner–Nordström 解一样，这里有三种可能性：$GM>a$、$GM=a$ 和 $GM<a$。最后一种情形的特征是存在裸奇点，而极端情形 $GM=a$ 是不稳定的，正如 Reissner–Nordström 情形一样。由于这些情形在物理上的重要性较低，我们将集中关注 $GM>a$。这时 $\Delta$ 在两个半径处为零，由
\begin{equation*}
r_{\pm} = GM \pm \sqrt{G^{2}M^{2}-a^{2}}.
\tag{6.82}
\end{equation*}
给出。这两个半径都是类光曲面，我们将会看到它们正是事件视界；Kerr 黑洞的侧视图见图 6.7。对这些曲面的分析与 Reissner–Nordström 情形极为类似；不难找到穿过视界延拓的坐标。

由于 Kerr 是稳态而非静态的，$r_{\pm}$ 处的事件视界并不是渐近时间平移 Killing 矢量 $K=\partial_t$ 的 Killing 视界。$K^{\mu}$ 的模由
\begin{equation*}
K^{\mu}K_{\mu} = -\frac{1}{\rho^{2}}\left(\Delta - a^{2}\sin^{2}\theta\right).
\tag{6.83}
\end{equation*}
给出。

% ---- 书页 264 / p279 ----

\begin{figure}[H]
\centering
\includegraphics[width=0.72\textwidth]{fig_6.7.pdf}
\caption*{图 6.7\quad Kerr 解周围的视界结构（侧视图）。事件视界是类光曲面，它标示出一条分界：越过它之后就不可能再返回空间的某个区域。与此不同，静止极限面除了与事件视界相切之处（两极）以外都是类时的；它代表越过之后便无法保持静止观者身份的地方。静止极限面与外侧事件视界之间的能层（ergosphere）是一个可以进入、也可以再离开、但不能在其中保持静止的区域。}
\end{figure}

在外侧事件视界处它并不为零；事实上，在 $r=r_{+}$（此处 $\Delta=0$）我们有
\begin{equation*}
K^{\mu}K_{\mu} = \frac{a^{2}}{\rho^{2}}\sin^{2}\theta \geq 0.
\tag{6.84}
\end{equation*}
因此，这个 Killing 矢量在外侧视界处已经是类空的了，只有在北极和南极（$\theta=0,\pi$）处它是类光的。$K^{\mu}K_{\mu}=0$ 的点的轨迹当然就是静止极限面，由
\begin{equation*}
(r-GM)^{2} = G^{2}M^{2} - a^{2}\cos^{2}\theta,
\tag{6.85}
\end{equation*}
给出，而外侧事件视界则由
\begin{equation*}
(r_{+}-GM)^{2} = G^{2}M^{2} - a^{2}.
\tag{6.86}
\end{equation*}
给出。于是，这两个曲面之间存在一个区域，称为\textbf{能层}（ergosphere）。在能层内部，你必须沿着黑洞转动的方向运动（即 $\phi$ 方向）；不过，你仍然可以朝事件视界运动，也可以远离它（要退出能层也没有任何麻烦）。显然，甚至在穿越视界之前，能层就是一个可能发生有趣事情的地方；详情后述。

在急着画共形图之前，我们需要弄清楚真正的曲率奇点的性质；在这个时空中，它并不位于 $r=0$ 处，而是位于 $\rho=0$ 处（曲率不变量 $R_{\rho\sigma\mu\nu}R^{\rho\sigma\mu\nu}$ 在那里发散）。由于 $\rho^{2}=r^{2}+a^{2}\cos^{2}\theta$ 是两个明显非负的量之和，只有当两者同时为零时它才可能为零，即
\begin{equation*}
r = 0,\qquad \theta = \frac{\pi}{2}.
\tag{6.87}
\end{equation*}

% ---- 书页 265 / p280 ----

这个结果看上去有点滑稽，但请记住 $r=0$ 不是空间中的一个点，而是一个圆盘；满足 $r=0$、$\theta=\pi/2$ 的点集实际上就是这个圆盘边缘上的\emph{环}。转动把 Schwarzschild 奇点“软化”了，将它摊开成一个环。

如果你走进环内会发生什么？仔细的解析延拓（我们不打算真的去做）将会揭示，你会到达另一个渐近平直的时空，但它并不是你来自的那个时空的全同副本。新时空由 $r<0$ 的 Kerr 度规描述。结果是 $\Delta$ 永不为零，从而不存在视界。共形图（图 6.8）与 Reissner–Nordström 的情形非常相像，只是现在你可以穿过奇点。由于 Kerr 度规不是球对称的，共形图也就不像前面几种情形那样忠实；图上的一个点代表固定的 $t$ 和 $r$ 值，而对于不同的 $\theta$ 值，它将有不同的几何。

\begin{figure}[H]
\centering
\includegraphics[width=0.72\textwidth]{fig_6.8.pdf}
\caption*{图 6.8\quad $G^{2}M^{2}>a^{2}$ 的 Kerr 解的共形图。与类似的带电解一样，黑洞外部的区域存在着无穷多个副本。}
\end{figure}

% ---- 书页 266 / p281 ----

我们不仅有这些经由黑洞与我们所在区域相连的、彼此不同的渐近平直区域所带来的通常怪异性，而且环奇点附近的区域还有额外的病态之处：闭合类时曲线（closed timelike curves，CTC）。如果你考虑一些在 $\phi$ 方向绕行、同时保持 $\theta$ 和 $t$ 固定、$r$ 取很小负值的轨迹，沿这样一条路径的线元是
\begin{equation*}
ds^{2} \approx a^{2}\left(1+\frac{2GM}{r}\right)\mathrm{d}\phi^{2},
\tag{6.88}
\end{equation*}
对于很小的负 $r$，它是负的。由于这些路径是闭合的，它们显然就是 CTC。因此，你有可能在过去遇见你自己，以及随之而来的一切。

当然，我们关于 Kerr 解析延拓所说的一切，都带有我们在讨论 Schwarzschild 和 Reissner–Nordström 时提到过的同样的告诫；真实的引力坍缩不太可能形成这些古怪的时空。尽管如此，拥有精确解总是有用的。此外，对于 Kerr 度规，即使我们停留在事件视界之外，奇怪的事情也在发生，现在我们就转向这些事情。

我们首先更仔细地考察黑洞的角速度。显然，常规的角速度定义必须稍加修改，才能应用于像时空度规这样抽象的对象。让我们考虑在 Kerr 黑洞的赤道面（$\theta=\pi/2$）内、半径 $r$ 处沿 $\phi$ 方向发射的一个光子的命运。它发射的那一瞬间，其动量在 $r$ 和 $\theta$ 方向都没有分量，因此轨迹为类光的条件是
\begin{equation*}
ds^{2} = 0 = g_{tt}\mathrm{d}t^{2} + g_{t\phi}\left(\mathrm{d}t\,\mathrm{d}\phi + \mathrm{d}\phi\,\mathrm{d}t\right) + g_{\phi\phi}\mathrm{d}\phi^{2}.
\tag{6.89}
\end{equation*}
这可以立即解出
\begin{equation*}
\frac{\mathrm{d}\phi}{\mathrm{d}t} = -\frac{g_{t\phi}}{g_{\phi\phi}} \pm \sqrt{\left(\frac{g_{t\phi}}{g_{\phi\phi}}\right)^{2} - \frac{g_{tt}}{g_{\phi\phi}}}.
\tag{6.90}
\end{equation*}
如果在 Kerr 度规的静止极限面上计算这个量，就有 $g_{tt}=0$，两个解为
\begin{equation*}
\frac{\mathrm{d}\phi}{\mathrm{d}t} = 0,\qquad \frac{\mathrm{d}\phi}{\mathrm{d}t} = \frac{a}{2G^{2}M^{2}+a^{2}}.
\tag{6.91}
\end{equation*}
非零解与 $a$ 同号；我们把它解释为光子绕黑洞运动的方向与黑洞转动的方向相同。零解则意味着，沿逆着黑洞转动方向射出的光子在这个坐标系中完全不动。注意，我们并没有给出光子轨迹的完整解，只是证明了它的瞬时速度为零。这就是著名的“参考系拖曳”（dragging of inertial frames）现象的一个例子；
% ---- 书页 267 / p282 ----
它将在第 7 章的一道习题中得到更多探讨。有质量的粒子必定比光子运动得慢，一旦进入静止极限面之内，它们就不可避免地被拖着随黑洞一起转动。当我们向位于 $r_{+}$ 处的外侧事件视界靠近时，这种拖曳仍在继续；我们可以定义事件视界本身的角速度 $\Omega_{\mathrm{H}}$，即粒子在视界处的最小角速度。直接由式 (6.90) 得到
\begin{equation*}
\Omega_{\mathrm{H}} = \left(\frac{\mathrm{d}\phi}{\mathrm{d}t}\right)_{-}(r_{+}) = \frac{a}{r_{+}^{2}+a^{2}}.
\tag{6.92}
\end{equation*}

\section{Penrose 过程与黑洞热力学}

黑洞热力学是广义相对论中最引人入胜、也最神秘的课题之一。不过，为了走到那里，让我们先从一件看上去非常直截了当的事情开始：Kerr 度规中沿测地线的运动。我们知道，如果考虑与 Killing 矢量 $K=\partial_t$ 和 $R=\partial_\phi$ 相联系的守恒量，这样的讨论将会得到简化。就眼前的目的而言，我们可以把注意力限制在有质量的粒子上，对它们可以使用四维动量
\begin{equation*}
p^{\mu} = m\frac{\mathrm{d}x^{\mu}}{\mathrm{d}\tau},
\tag{6.93}
\end{equation*}
其中 $m$ 是粒子的静质量。于是，我们可以取粒子真实的能量和角动量作为我们的两个守恒量，
\begin{equation*}
E = -K_{\mu}p^{\mu} = m\left(1-\frac{2GMr}{\rho^{2}}\right)\frac{\mathrm{d}t}{\mathrm{d}\tau} + \frac{2mGMar}{\rho^{2}}\sin^{2}\theta\,\frac{\mathrm{d}\phi}{\mathrm{d}\tau}
\tag{6.94}
\end{equation*}
以及
\begin{equation*}
L = R_{\mu}p^{\mu} = -\frac{2mGMar}{\rho^{2}}\sin^{2}\theta\,\frac{\mathrm{d}t}{\mathrm{d}\tau} + \frac{m(r^{2}+a^{2})^{2}-m\Delta a^{2}\sin^{2}\theta}{\rho^{2}}\sin^{2}\theta\,\frac{\mathrm{d}\phi}{\mathrm{d}\tau}.
\tag{6.95}
\end{equation*}
这些定义与上一章中使用的守恒量定义有所不同，在上一章里 $E$ 和 $L$ 取的是\emph{每单位质量}的能量和角动量。当然，无论哪种取法，它们都是守恒的。

$E$ 的定义中有一个负号，这是因为在无穷远处 $K^{\mu}$ 和 $p^{\mu}$ 都是类时的，它们的内积为负，而我们希望能量是正的。然而在能层内部，$K^{\mu}$ 变成了类空的；因此我们可以设想有这样的粒子，使得
\begin{equation*}
E = -K_{\mu}p^{\mu} < 0.
\tag{6.96}
\end{equation*}

% ---- 书页 268 / p283 ----

这件事令我们感到的不安，多少可以因如下认识而缓解：只要在静止极限面之外，\emph{所有}粒子的能量都必定为正；因此，能层内部具有负能量的粒子，要么留在能层之中，要么若想逃出去，就必须被加速到能量变为正。

尽管如此，这一认识仍然开辟了一条从转动黑洞中汲取能量的途径；这个方法就是\textbf{Penrose 过程}（Penrose process）。想法很简单：从能层之外出发，你给自己装备一块大石头，跃向黑洞。如果我们把（你 $+$ 石头）这个系统的四维动量记为 $p^{(0)\mu}$，那么能量 $E^{(0)}=-K_{\mu}p^{(0)\mu}$ 无疑是正的，而且当你沿自己的测地线运动时它是守恒的。一旦你进入能层，就用尽全力、以非常特别的方式把石头猛掷出去。如果把你的动量记为 $p^{(1)\mu}$、石头的动量记为 $p^{(2)\mu}$，那么在你掷出石头的那一瞬间，正如狭义相对论中一样，我们有动量守恒：
\begin{equation*}
p^{(0)\mu} = p^{(1)\mu} + p^{(2)\mu}.
\tag{6.97}
\end{equation*}
与 Killing 矢量 $K_{\mu}$ 缩并，给出
\begin{equation*}
E^{(0)} = E^{(1)} + E^{(2)}.
\tag{6.98}
\end{equation*}
但是，如果我们设想你足够强壮（而且投掷得足够精准），你就可以按照式 (6.96) 来安排你的投掷，使得 $E^{(2)}<0$。此外，Penrose 还能够证明，你可以像图 6.9 那样安排初始轨迹和投掷，使得随后你沿一条测地线轨迹回到静止极限面之外、进入外部宇宙。由于你的能量沿途守恒，

\begin{figure}[H]
\centering
\includegraphics[width=0.72\textwidth]{fig_6.9.pdf}
\caption*{图 6.9\quad Penrose 过程（俯视图）。一个物体落向 Kerr 黑洞，并在能层之内（静止极限面以内、外侧事件视界之外）分裂为二。一块以负能量 $E^{(2)}$ 落入视界，另一块则以比原来落入物体更大的能量逃逸到无穷远。}
\end{figure}

% ---- 书页 269 / p284 ----

最终我们将有
\begin{equation*}
E^{(1)} > E^{(0)}.
\tag{6.99}
\end{equation*}
这样，你脱身而出时携带的能量就比你进入时\emph{更多}。

世上没有免费的午餐；你获得的能量来自某处，而那个某处就是黑洞。事实上，Penrose 过程是通过减小转动黑洞的角动量来从中汲取能量的；要让这个戏法奏效，你必须逆着黑洞的转动方向掷出石头。为了更精确地看清这一点，回忆一下我们在本章前文中的论断：稳态时空中的任何事件视界，对于某个 Killing 矢量来说都是一个 Killing 视界。对 Kerr 而言，这个矢量是时间平移和转动 Killing 矢量的线性组合，
\begin{equation*}
\chi^{\mu} = K^{\mu} + \Omega_{\mathrm{H}}R^{\mu},
\tag{6.100}
\end{equation*}
其中 $\Omega_{\mathrm{H}}$ 恰恰就是式 (6.92) 中定义的视界角速度。利用 $K=\partial_t$ 和 $R=\partial_\phi$，不难验证 $\chi^{\mu}$ 在外侧事件视界处变为类光。具有动量 $p^{(2)\mu}$ 的粒子“在时间中向前”地穿越事件视界，这一陈述不过是
\begin{equation*}
p^{(2)\mu}\chi_{\mu} < 0.
\tag{6.101}
\end{equation*}
代入 $E$ 和 $L$ 的定义，我们看到这个条件等价于
\begin{equation*}
L^{(2)} < \frac{E^{(2)}}{\Omega_{\mathrm{H}}}.
\tag{6.102}
\end{equation*}
既然我们已经安排 $E^{(2)}$ 为负、$\Omega_{\mathrm{H}}$ 为正，这就意味着粒子必须带有负的角动量——它在逆着黑洞转动的方向运动。一旦你逃出能层、石头落入事件视界之内，黑洞的质量和角动量就是原来的数值再加上石头所带来的负的贡献：
\begin{gather*}
\delta M = E^{(2)}\\
\delta J = L^{(2)},
\tag{6.103}
\end{gather*}
其中 $J=Ma$ 是黑洞的角动量。于是式 (6.102) 就变成了对角动量能够减小多少的一个限制：
\begin{equation*}
\delta J < \frac{\delta M}{\Omega_{\mathrm{H}}}.
\tag{6.104}
\end{equation*}
如果我们恰好达到这个极限——随着我们掷入的石头越来越接近类光——我们就得到了“理想”过程，其中 $\delta J = \delta M/\Omega_{\mathrm{H}}$。

现在我们将利用这些想法来验证：尽管你可以用 Penrose 过程从黑洞中汲取能量（从而减小 $M$），你却不能

% ---- 书页 270 / p285 ----

违反面积定理：事件视界的面积是不减的。虽然质量减小了，但角动量也必须减小，而且两者减小的组合方式只允许面积增加。为了看清这一点，我们来计算外侧事件视界的面积，它位于
\begin{equation*}
r_{+} = GM + \sqrt{G^{2}M^{2}-a^{2}}.
\tag{6.105}
\end{equation*}
视界上的诱导度规 $\gamma_{ij}$（其中 $i$ 和 $j$ 取遍 $\{\theta,\phi\}$）可以直截了当地求得：在式 (6.70) 中令 $r=r_{+}$（于是 $\Delta=0$）、$\mathrm{d}t=0$ 和 $\mathrm{d}r=0$：
\begin{align*}
\gamma_{ij}\,\mathrm{d}x^{i}\mathrm{d}x^{j} &= ds^{2}(\mathrm{d}t=0,\,\mathrm{d}r=0,\,r=r_{+})\\
&= (r_{+}^{2}+a^{2}\cos^{2}\theta)\mathrm{d}\theta^{2} + \left[\frac{(r_{+}^{2}+a^{2})^{2}\sin^{2}\theta}{r_{+}^{2}+a^{2}\cos^{2}\theta}\right]\mathrm{d}\phi^{2}.
\tag{6.106}
\end{align*}
于是，视界面积就是诱导体积元的积分，
\begin{equation*}
A = \int\sqrt{|\gamma|}\,\mathrm{d}\theta\,\mathrm{d}\phi.
\tag{6.107}
\end{equation*}
行列式为
\begin{equation*}
|\gamma| = (r_{+}^{2}+a^{2})^{2}\sin^{2}\theta,
\tag{6.108}
\end{equation*}
所以视界面积就是
\begin{equation*}
A = 4\pi(r_{+}^{2}+a^{2}).
\tag{6.109}
\end{equation*}

为了证明面积不会减小，更方便的做法是改用黑洞的\textbf{不可约质量}（irreducible mass），其定义为
\begin{align*}
M_{\mathrm{irr}}^{2} &= \frac{A}{16\pi G^{2}}\\
&= \frac{1}{4G^{2}}(r_{+}^{2}+a^{2})\\
&= \frac{1}{2}\left(M^{2}+\sqrt{M^{4}-(Ma/G)^{2}}\right)\\
&= \frac{1}{2}\left(M^{2}+\sqrt{M^{4}-(J/G)^{2}}\right).
\tag{6.110}
\end{align*}
经过一点计算，我们可以通过求导得到 $M_{\mathrm{irr}}$ 如何受质量或角动量变化的影响：
\begin{equation*}
\delta M_{\mathrm{irr}} = \frac{a}{4GM_{\mathrm{irr}}\sqrt{G^{2}M^{2}-a^{2}}}\left(\Omega_{\mathrm{H}}^{-1}\delta M - \delta J\right).
\tag{6.111}
\end{equation*}

% ---- 书页 271 / p286 ----

于是，我们的极限 (6.104) 就变成
\begin{equation*}
\delta M_{\mathrm{irr}} > 0.
\tag{6.112}
\end{equation*}
不可约质量永远无法减小；名称由此而来。由此可知，在把黑洞的转动减缓到零之前，我们能够从黑洞汲取的最大能量为
\begin{equation*}
M - M_{\mathrm{irr}} = M - \frac{1}{\sqrt{2}}\left(M^{2}+\sqrt{M^{4}-(J/G)^{2}}\right)^{1/2}.
\tag{6.113}
\end{equation*}
这种彻底汲取的结果，是一个质量为 $M_{\mathrm{irr}}$ 的 Schwarzschild 黑洞。事实证明，我们所能做到的最好情形是从一个极端 Kerr 黑洞出发；这样我们可以取出它总能量的大约 29\%。

$M_{\mathrm{irr}}$ 的不可约性直接引出这样一个事实：面积 $A$ 永不减小。由式 (6.110) 和 (6.111) 我们有
\begin{equation*}
\delta A = 8\pi G\,\frac{a}{\Omega_{\mathrm{H}}\sqrt{G^{2}M^{2}-a^{2}}}\left(\delta M - \Omega_{\mathrm{H}}\delta J\right),
\tag{6.114}
\end{equation*}
它可以改写为
\begin{equation*}
\delta M = \frac{\kappa}{8\pi G}\delta A + \Omega_{\mathrm{H}}\delta J,
\tag{6.115}
\end{equation*}
其中我们引入了
\begin{equation*}
\kappa = \frac{\sqrt{G^{2}M^{2}-a^{2}}}{2GM\left(GM+\sqrt{G^{2}M^{2}-a^{2}}\right)}.
\tag{6.116}
\end{equation*}
量 $\kappa$ 当然就是 Kerr 解的表面引力，你可以把式 (6.100) 代入式 (6.9) 来验证这一点。

像 (6.115) 这样的方程，最先促使人们开始思考黑洞与热力学之间的对应。考虑热力学第一定律，
\begin{equation*}
\mathrm{d}E = T\,\mathrm{d}S - p\,\mathrm{d}V,
\tag{6.117}
\end{equation*}
其中 $T$ 是温度，$S$ 是熵，$p$ 是压强，$V$ 是体积，因此 $p\,\mathrm{d}V$ 项代表我们对系统所做的功。很自然地，我们可以把式 (6.115) 中的 $\Omega_{\mathrm{H}}\delta J$ 项看成我们向黑洞投掷石头时对它所做的功。这样，如果我们把热力学量——能量、熵和温度——与黑洞的质量、面积和表面引力等同起来，这个对应就开始成形：
\begin{gather*}
E \leftrightarrow M\\
S \leftrightarrow A/4G\\
T \leftrightarrow \kappa/2\pi.
\tag{6.118}
\end{gather*}

% ---- 书页 272 / p287 ----

（记住，我们使用的单位制中 $\hbar=c=k=1$。）在经典广义相对论的语境中，这个类比本质上是完美的：热力学的每一条定律都对应着一条黑洞力学定律。处于热平衡的系统将弛豫到一个稳态，对应于稳态黑洞。热力学第零定律说，在热平衡中温度在整个系统中处处相同；对黑洞的类似陈述是：稳态黑洞在整个视界上具有恒定的表面引力。至少在使事件视界成为 Killing 视界的那同样的合理假设之下，这一点将会成立。正如我们已经看到的，第一定律 (6.117) 等价于 (6.115)。第二定律——熵永不减小——不过就是视界面积永不减小这一陈述。最后，第三定律的通常表述是：在任何物理过程中都不可能达到 $T=0$，或者说，当温度趋于零时熵必须趋于零。对黑洞来说，这不太灵验；事实证明 $\kappa=0$ 对应于极端黑洞，而极端黑洞的面积不一定为零。不过，热力学第三定律本身其实也不真正灵验，因为确实存在违反它的普通物理系统；第三定律适用于某些情形，但并不是真正基本的定律。

在提出对应关系 (6.118) 时，我们有点儿作弊了；你会注意到，把 $T\mathrm{d}S$ 与 $\kappa\mathrm{d}A/8\pi G$ 等同起来之后，我们并不知道如何分别归一化 $S/A$ 或 $T/\kappa$，只知道它们的组合。不过，正如我们将在第 9 章讨论的，Hawking 证明了黑洞背景中的量子场使得黑洞以温度 $T=\kappa/2\pi$ 发出辐射。一旦知道了这一点，我们就可以把 $A/4G$ 解释为黑洞的真实熵。Bekenstein 提出了\textbf{广义第二定律}（generalized second law）：物质与黑洞合起来的熵永不减小：
\begin{equation*}
\delta\left(S + \frac{A}{4G}\right) \geq 0.
\tag{6.119}
\end{equation*}
广义第二定律实际上可以在多种不同的假设下得到证明。不过，我们通常喜欢把系统的熵与可到达的量子态数目的对数联系起来。因此，这个概念与无毛定理之间存在着某种张力：无毛定理告诉我们，具有固定电荷、质量和自旋的黑洞，其可能的状态极少（事实上只有一个）。这种行为似乎在向我们暗示，它是量子力学与引力之间相互作用之深刻性质的一个表征。

\section{习题}

\textbf{1.}\quad 证明：耦合的 Einstein–Maxwell 方程组可以同时由度规 (6.62) 和静电势 (6.67) 求解，只要 $H(\vec{x})$ 满足拉普拉斯方程
\begin{equation*}
\nabla^{2}H = 0.
\tag{6.120}
\end{equation*}

\textbf{2.}\quad 考虑 Kerr 黑洞赤道面内无质量粒子的轨道，仿射参量为 $\lambda$。

\textbf{(a)}\quad 证明
\begin{equation*}
\left(\frac{\mathrm{d}r}{\mathrm{d}\lambda}\right)^{2} = \frac{\Sigma^{2}}{\rho^{4}}\left(E - LW_{+}(r)\right)\left(E - LW_{-}(r)\right),
\tag{6.121}
\end{equation*}
其中 $\Sigma^{2} = (r^{2}+a^{2})^{2} - a^{2}\Delta(r)\sin^{2}\theta$，$E$ 和 $L$ 是守恒的能量和角动量，而 $W_{\pm}(r)$ 的表达式则要你自己求出。

\textbf{(b)}\quad 利用这个结果，并假设 $\Sigma^{2}$ 处处大于零，证明：赤道面内光子的轨道不可能在外侧事件视界 $r_{+}$ 之内存在转折点。这意味着，入射光线一旦穿越 $r_{+}$ 就无法逃逸，所以它确实是一个事件视界。

% ---- 书页 273 / p288 ----

\textbf{3.}\quad 在存在电磁场的情况下，电荷为 $e$、质量为 $m$ 的粒子服从
\begin{equation*}
\frac{\mathrm{d}^{2}x^{\mu}}{\mathrm{d}\tau^{2}} + \Gamma^{\mu}_{\rho\sigma}\frac{\mathrm{d}x^{\rho}}{\mathrm{d}\tau}\frac{\mathrm{d}x^{\sigma}}{\mathrm{d}\tau} = \frac{e}{m}F^{\mu}{}_{\nu}\frac{\mathrm{d}x^{\nu}}{\mathrm{d}\tau}.
\tag{6.122}
\end{equation*}
设想这样一个粒子在电荷为 $Q$、质量为 $M$ 的 Reissner–Nordström 黑洞的场中运动。

\textbf{(a)}\quad 证明能量
\begin{equation*}
E = m\left(1 - \frac{2GM}{r} + \frac{GQ^{2}}{r^{2}}\right)\frac{\mathrm{d}t}{\mathrm{d}\tau} + \frac{eQ}{r}
\tag{6.123}
\end{equation*}
是守恒的。

\textbf{(b)}\quad Penrose 型的过程对带电黑洞行得通吗？对于最大的物理过程，黑洞质量的改变量 $\delta M$ 是多少？

\textbf{4.}\quad 考虑静态坐标下的 de Sitter 空间：
\begin{equation*}
ds^{2} = -\left(1 - \frac{\Lambda}{3}r^{2}\right)\mathrm{d}t^{2} + \frac{\mathrm{d}r^{2}}{1 - \frac{\Lambda}{3}r^{2}} + r^{2}\mathrm{d}\Omega^{2}.
\end{equation*}
这个空间有一个 Killing 矢量 $\partial_t$，它在 $r=0$ 附近是类时的，而在某个 Killing 视界上是类光的。试定出 Killing 视界的径向位置 $r_{\mathrm{K}}$。该视界的表面引力 $\kappa$ 是多少？考虑通过替换 $t\rightarrow i\tau$ 得到的 de Sitter 空间的欧氏号差版本。证明：只要把 $\tau$ 取成周期的，就可以通过一个坐标变换使欧氏度规在视界处正则。

\textbf{5.}\quad 一位观者沿周长为 $2\pi R$ 的圆轨道，绕着电荷为 $Q$、质量为 $M$ 的 Reissner–Nordström 黑洞运行，他所看到的磁场是什么？

\textbf{6.}\quad 考虑一个 Kerr 黑洞，在其赤道面内有一个质量可忽略的吸积盘（accretion disk）。假设盘中的粒子沿测地线运动（也就是说，忽略任何压强支撑）。现在进一步假设盘中含有一些铁原子，它们正被某个辐射源激发。当铁原子退激发时，它们会发出辐射，在它们自己的静止系中测量，其频率为已知的 $\nu_{0}$。假设我们在远离黑洞的地方探测这种辐射（我们也处在赤道面内）。从盘的任一边缘以及从盘的中心发出的光子，其观测频率各是多少？考虑盘与黑洞同向转动和反向转动这两种情形。我们能否利用这些测量来确定黑洞的质量和角动量？
